\section{Trading \& Markets}

\begin{frame}{Trading Introduction}
  We will introduce the concepts of
  \begin{itemize}
    \item Trading
    \item Exchanges
    \item Orders and the orderbook
  \end{itemize}

  \vspace{0.3cm}

  \textcolor{sqt-question}{
    I thought that firms don't care about prior knowledge of finance?
  }

  \vspace{0.3cm}

  \textbf{Motivation:}
  \begin{itemize}
    \item this is only true during the interview stage for large HFT firms
    \item beginner understanding of orderbook dynamics is useful for
    \begin{itemize}
      \item mock trading games
      \item general understanding of the industry
    \end{itemize}
  \end{itemize}
\end{frame}

\begin{frame}{Spot Trades}
  \begin{sqt:tip}
    Alice says to Bob:
    \begin{center}
      ``I want to to sell you 5 AAPL shares. Each share will be sold at \$30.
      This transfer will occur by the end of the next business day (tomorrow)."
    \end{center}
  \end{sqt:tip}

  This defines a \textbf{spot trade} (as opposed to a \textbf{future trade}).

  \begin{sqt:definition}
    A \textbf{spot trade} is an agreement to exchange assets for money, with
    the actual transfer occuring almost immediately (within 1 or 2 business days).
    \begin{itemize}
      \item sides (Alice is selling, Bob is buying)
      \item quantity of the asset being traded ($5$)
      \item price per asset ($\$30$)
    \end{itemize}
  \end{sqt:definition}
\end{frame}

\begin{frame}{Trade Execution and Settlement}
  Once Bob \textbf{agrees} to this trade, the trade can be \textbf{executed}.

  \begin{sqt:definition}
    Upon trade \textbf{execution}, both parties are given \textbf{settlement obligations}.
  \end{sqt:definition}
  \begin{sqt:tip}
    For example, in Alice \& Bob's equity spot trade, the obligations are:
    \begin{itemize}
      \item Alice must give Bob $5$ shares of AAPL by the settlement deadline
      \item Bob must give Alice $5\times \$30 = \$150$ by the settlement deadline
    \end{itemize}
    Once the obligations are fulfilled, the trade is considered \textbf{settled}.
  \end{sqt:tip}

  Obligations on other asset classes and types of trades (e.g. futures and options)
  are considerably more complex than this.
  Concepts related to \textbf{clearing} and \textbf{settlement} are
  considered part of the \textbf{post-trade life-cycle}.
\end{frame}

\begin{frame}{Holdings and Position}
  At each point in time, each participant has a \textbf{holdings} and a \textbf{position} in each asset.
  \begin{sqt:definition}
    \textbf{Position} includes the effects of executed \textbf{trades} that have not been \textbf{settled} yet.
  \end{sqt:definition}

  \begin{center}
    \input{assets/holdings_vs_positions.tex}
  \end{center}
\end{frame}

\begin{frame}{Holdings and Position}
  \begin{sqt:result}
    Your \textbf{exposure} and \textbf{PnL} is determined by your \textbf{position}.
  \end{sqt:result}

  This is because the terms of the trade (price, quantity) are locked in immediately upon trade execution,
  so the transfer of assets and money can be thought of as occuring immediately.

  \begin{sqt:tip}
    For example, you agree to buy a share of AAPL at \$30.
    Immediately, you are ``exposed'' to the movement of AAPL stock price.
    If the market price goes down to \$10, you are still paying \$30 for it upon settlement,
    but might end up having to sell it away for \$10.
  \end{sqt:tip}

  As a result, for trading games and simulations, we can abstract away the entire post-trade life-cycle
  and just consider our asset and cash \textbf{position}, disregarding \textbf{holdings} entirely.
\end{frame}

\begin{frame}{What is an Exchange}
  \begin{sqt:definition}
    An \textbf{exchange} is a centralised venue for participants to\\
    find willing counterparties to make \textbf{trades} with.
  \end{sqt:definition}
  \begin{sqt:definition}
    If a trade does not occur on an exchange, then the trade is \textbf{over-the-counter} (OTC).
  \end{sqt:definition}
  \vspace{0.5cm}

  \begin{center}
    \begin{columns}
      \begin{column}[t]{0.3\textwidth}
        US exchanges
        \begin{itemize}
          \item (XNYS) NYSE
          \item (XNAS) Nasdaq
          \item (XCBO) CBOE
          \item (XCME) CME
        \end{itemize}
      \end{column}

      \begin{column}[t]{0.3\textwidth}
        APAC exchanges
        \begin{itemize}
          \item (XHKG) HKEX
          \item (XSHG) SSE
          \item (XSGE) SHFE
        \end{itemize}
      \end{column}
    \end{columns}
  \end{center}
\end{frame}

\begin{frame}{Limit Orders}
  Each participant maintains a set of \textbf{limit orders}, which the exchange uses
  to automatically infer whether they would agree to any hypothetical trade or not.\\
  Here is what a limit order may look like:
  \begin{sqt:tip}
    Alice tells the exchange
    \begin{center}
      ``I am willing to sell a maximum of 5 AAPL shares.\\
      I am willing to sell each share at a price of \$30 or better.''
    \end{center}
  \end{sqt:tip}

  This is called \textbf{inserting} an order.

  \begin{center}
    \input{assets/limit_order_intro.tex}
  \end{center}
\end{frame}

\begin{frame}{Bids and Offers}
  By submitting the \textbf{sell order}, Alice authorizes the exchange to accept trades on her behalf,
  until 5 AAPL shares have been \textbf{sold}. Once this happens, the order is considered \textbf{filled}.

  \begin{sqt:definition}
    \begin{itemize}
      \item A \textbf{bid} (buy order) expresses a willingness to buy
      \item An \textbf{offer/ask} (sell order) expresses a willingness to sell
    \end{itemize}
  \end{sqt:definition}

  \begin{center}
    \input{assets/bids_and_asks.tex}
  \end{center}
\end{frame}

\begin{frame}{Orderbook}
  \begin{sqt:definition}
    The exchange keeps track of each participant's \textbf{open orders} (unfilled),\\
    and collects them into an \textbf{orderbook}.
  \end{sqt:definition}

  \begin{center}
    \input{assets/orderbook.tex}
  \end{center}
\end{frame}

\begin{frame}{Limit Orders In Cross}

  \begin{sqt:definition}
    Two limit orders are \textbf{in cross} if
    \begin{itemize}
      \item One order is a \textbf{bid} and the other is an \textbf{ask}.
      \item There exists an \textbf{intersection} between the price ranges they are willing to trade at
    \end{itemize}
  \end{sqt:definition}
  \begin{center}
    \input{assets/in_cross.tex}
  \end{center}
\end{frame}
